\begin{frame}{자주 하는 실수: 절댓값을 그대로 해석하기}
  \begin{itemize}
    \item 학습된 $(7.264,-3.643)$이 진짜 $(2,-1)$과 스케일이
      다르다 $\to$ ``모델이 부정확하다''로 오해하기 쉬움.
    \item Bradley-Terry 손실은 오직 \textbf{상대 차이}만 학습
      신호로 쓰므로, \textbf{절대 스케일은 원래 정해지지 않음}.
    \item 실전에서 보상모델 출력을 ``만족도 점수''처럼 절대적으로
      해석하면 안 됨 — 항상 \textbf{다른 후보와 비교했을 때}
      상대적으로 얼마나 나은가로만 읽어야.
  \end{itemize}
  \vskip0.6em
  \begin{block}{FAQ: 300개는 숫자적으로 특별한가 / tie 처리}
    쌍 수보다 \textbf{쌍이 특징 공간의 서로 다른 방향을 얼마나
    대표하느냐}가 핵심 — 30개면 방향이 고르게 안 커버돼 비율이
    $-2$에서 꽤 벗어날 수 있음. ``둘 다 비슷하다''(tie)를
    win/lose에 억지로 넣으면 \textbf{잡음} — 실전에선 걸러내거나,
    연속 선호(0$\sim$1 실수값)로 확장.
  \end{block}
\end{frame}
